\begin{abstract}
Industrial fabric inspection requires detecting tiny and weakly visible defects
under constrained computation. Low-resolution detectors are efficient but miss
defects whose visual evidence is insufficient, whereas full high-resolution
inference is expensive and may even degrade easy cases due to scale-distribution
shift. This paper proposes a visibility-boundary-aware detection framework for
industrial fabric defect inspection. Through a cross-resolution failure
diagnosis we show that a range of cheap interventions---threshold calibration,
loss reweighting, knowledge distillation, and sparse patch reinspection---fail
to recover visibility-limited misses, because the missing evidence is not
present in the low-resolution input. Guided by this diagnosis, the proposed
framework uses cross-resolution failure supervision to learn which images are
visibility-limited and selectively triggers full-image high-resolution
reinspection only on those, while explicitly separating ambiguous defects that
remain unrecoverable even at high resolution. On a public fabric defect
detection dataset from Science Data Bank, the method recovers about $69\%$ of
the high-resolution recall gain at roughly $63\%$ of its computational cost, and
an oracle triage even exceeds uniform high-resolution inference by retaining
low-resolution inference on easy images. The analysis further reveals that a
large fraction of missed defects are inherently ambiguous or
annotation-sensitive, highlighting practical reliability boundaries for
AI-based fabric inspection.
\end{abstract}

\begin{keyword}
fabric defect detection \sep industrial visual inspection \sep
adaptive-resolution inference \sep cross-resolution failure analysis \sep
reliability-aware detection
\end{keyword}
